\begin{abstract}
AI-assisted software development is no longer a matter of isolated prompts. It has reorganized itself around frameworks that structure the work with persistent artifacts, agentic roles, integrated execution, and validation. What has been missing is a review that takes these operational frameworks as its object, through a software engineering lens, rather than treating agent architectures in the abstract or cataloguing techniques tied to single lifecycle tasks. This article presents a systematic literature review conducted under the Kitchenham protocol. The review gathered 37 studies published between 2022 and 2026 after two screening passes and three rounds of snowballing, including only academic papers and archived preprints with verifiable records. A hybrid thematic synthesis derived six inductive themes and tested them against a proposed taxonomy of six dimensions: specification, context, roles, execution, validation, and portability. The evidence supports the shift from prompt to process, which holds in 34 of the 37 studies, and confirms five of the six dimensions. Portability is the weak one, surfacing almost always as a diagnostic risk rather than a design property, and two axes emerge outside the taxonomy: human-AI collaboration and security. We identify three recurring axes of tension and show that context and traceability, not code generation itself, are the field's transversal bottlenecks. We close with a process-oriented research agenda built around benchmarks that evaluate the pipeline, head-to-head comparison between frameworks, and longitudinal evaluation of governance.
\end{abstract}

\noindent\textbf{Keywords:} AI-assisted software development; AI agents; systematic review; software engineering; agentic frameworks.
